% Standalone source. This investigation does not compile the manuscript.
% Established source results and direct derived comparisons are distinguished.
\documentclass[11pt]{article}
\usepackage[T1]{fontenc}
\usepackage[utf8]{inputenc}
\usepackage[margin=1in]{geometry}
\usepackage{amsmath,amssymb,amsthm}
\usepackage[hidelinks]{hyperref}
\newtheorem{theorem}{Theorem}[section]
\newtheorem{proposition}[theorem]{Proposition}
\newtheorem{lemma}[theorem]{Lemma}
\theoremstyle{definition}\newtheorem{definition}[theorem]{Definition}
\theoremstyle{remark}\newtheorem{remark}[theorem]{Remark}
\newtheorem{example}[theorem]{Example}
\newcommand{\Set}{\mathbf{Set}}
\newcommand{\id}{\mathrm{id}}
\newcommand{\Ctx}{\operatorname{Ctx}}
\newcommand{\im}{\operatorname{im}}
\title{Changing operations and observations through contextad comparisons}
\author{A source-level comparison and its exact obligations}
\date{5 November 2026}
\begin{document}\maketitle

\section{The useful connection}
The earlier contextad construction can supply a compositional source for
changes of operational language. It does not itself choose the observations
that determine a characterization. There are two logically separate pieces:
a morphism of operational structures, and a backwards translation of admitted
tests that exactly represents their responses. With both pieces, the existing
Chu comparison transports response quotients, concept completions and partial
Dowker maps coherently. A strong operational comparison alone supplies none of
the requisite observational adequacy or native reconstruction conclusions.

This note gives the typed connection, a reader example, and a display-span
obstruction. The results about contextad morphisms are existing mathematics;
the characterizations reuse the established bridge proved in the earlier
common-generation manuscript. Their assembly is a direct derivation, without a
claim of historical novelty.

\subsection{A necessary source update}
The restored manuscripts dated 1 November used Capucci--Myers
\cite{CM1}, the 84-page version titled \emph{Contextads as Wreaths; Kleisli,
Para, and Span Constructions as Wreath Products}. Version 2, dated 3 November
2026, is now available as the 120-page \emph{Contextful Arrows via Wreath
Products} \cite{CM2}. In v2, Definition 5.25 specifies morphisms of contextads;
Examples 5.33--5.35 identify comonad, action and display-map instances;
Definition 5.36 and equations (5.6.20)--(5.6.21) give the induced strong double
functor. These are the precise antecedents used here.

This version update does not remove every foundational qualification.
Section 5.4 explicitly conditions its concrete presentation of the free
Kleisli tricocompletion on a conjectural identification, and Remark 5.20 takes
the higher wreath-product formula for granted. The concrete computations
below can be checked directly from the displayed structure maps and do not
require declaring those qualifications settled. In particular, v1 theorem
numbers must not be attached to the changed v2 title without a version label.

\section{The operational map and its variance}
Let a contextad have indexed category $K:C^{\mathrm{op}}\to\mathbf{Cat}$,
extension $e:\int K\to C$, unit grade $I_A$, and accumulation
$P\otimes Q\in K_A$ for $P\in K_A$, $Q\in K_{e_AP}$.
Its contextual arrows are
\[
 (P,f):A\Rightarrow B,\qquad f:e_AP\longrightarrow B.
\]
Writing $e_f(R):e_{e_AP}(f^*R)\to e_BR$ for extension of the cartesian lift,
composition has grade $P\otimes f^*R$ and body
\[
 e_A(P\otimes f^*R)\xrightarrow{\delta}e_{e_AP}(f^*R)
 \xrightarrow{e_f(R)}e_BR\xrightarrow{g}D.
\]
The pullback $f^*R$ changes a requirement along one actual body \emph{within}
one operational doctrine. It is not, by itself, a change of the whole doctrine.

For a second contextad $(K',e')$ on $C'$, a strong contextad morphism as in
\cite[Definition 5.25]{CM2} consists of a functor $F:C\to C'$, a cartesian
functor $G:\int K\to\int K'$ over $F$, and a natural lineator
\begin{equation}\label{eq:lineator}
 \ell_{A,P}:e'_{FA}(G_AP)\longrightarrow F(e_AP).
\end{equation}
Here ``cartesian'' means preservation of cartesian lifts, not that the square
of total and base categories is asserted to be a pullback square. It gives
coherent isomorphisms
\[
 G_A(u^*P')\cong(Fu)^*(G_{A'}P')\qquad(u:A\to A').
\]
There are also coherent invertible grade comparisons
\begin{align*}
 G_AI_A&\cong I'_{FA},\\
 G_A(P\otimes Q)&\cong
 G_AP\otimes'\ell_{A,P}^{*}(G_{e_AP}Q),
\end{align*}
compatible with $\varepsilon,\delta$, including the unit and associativity
coherences. The lineator \eqref{eq:lineator} need not be invertible. This is
not a contradiction: the strength refers to the unit and multiplication
comparisons, not to the lineator.

\begin{proposition}[The existing operational transport, explicitly typed]
Under these hypotheses the assignment
\begin{equation}\label{eq:ctx-map}
 T(P,f)=(G_AP,F(f)\ell_{A,P})
\end{equation}
defines the corresponding strong double comparison. On tight boundaries it
acts by $F$ and on grade cells by $G$, with the cartesian identification where
necessary. The map preserves contextual composition up to the specified
invertible comparison.
\end{proposition}
\begin{proof}
The body in \eqref{eq:ctx-map} starts at $e'_{FA}(G_AP)$ and ends at $FB$.
For successive $(P,f)$ and $(R,g)$, the image of the source composite has grade
$G_A(P\otimes f^*R)$. Its prescribed multiplication comparison and the
cartesian property identify this with
\[
 G_AP\otimes'\ell_{A,P}^{*}G_{e_AP}(f^*R)
 \cong G_AP\otimes'(F(f)\ell_{A,P})^*(G_BR),
\]
which is exactly the grade of the composite of the two images. The
multiplication compatibility of $\ell$ identifies the bodies. The unit law
identifies the image of $(I_A,\varepsilon_A)$ with the target unit. Their
coherence is exactly the assumed triangle and pentagon coherence of the grade
comparisons. A square with grade map $\alpha:P\to u^*P'$ and body equation
$v f=f'e_u(P')e_A\alpha$ is sent to the corresponding square by naturality
of $\ell$ and the cartesian comparison. Thus all endpoint squares, not only
globular cells, are transported. This is the concrete calculation underlying
\cite[equations (5.6.20)--(5.6.21)]{CM2}; it is not a new construction.
\end{proof}

The relative-construction manuscript \cite{Relative} already permitted a
base-changing strong operational comparison $T:B\to B'$ together with an
invertible pure-compatibility icon $TJ\Rightarrow J'F$. It also specified
lax double comparisons separately. The contextad morphism above supplies one
actual presentation-level origin for the strong class. It does not identify
all competing lax, oplax and pseudo morphism conventions.

\section{The additional observational datum}
An operational double category is not automatically a deterministic state
system. Supply an actual interpreted family of state sets $X_a$, admitted
operations $u:X_a\to X_b$, and Boolean tests $E_a$. They may come from a
contextad-based grammar, but this interpretation and the admissible tests
must be provided. In particular, neither an arbitrary noninvertible native
cell nor an arbitrary stochastic arrow is silently interpreted as equality of
deterministic states.

A type-changing operational comparison supplies a translation $a\mapsto Fa$
and interpreted maps $f_a:X_a\to Y_{Fa}$. Require
\begin{equation}\label{eq:state-nat}
 f_bu=\widetilde u f_a,
\end{equation}
where $\widetilde u:Y_{Fa}\to Y_{Fb}$ is the translated operation.
Write $R_a\subseteq X_a\times E_a$ and
$S_{Fa}\subseteq Y_{Fa}\times D_{Fa}$ for response incidence.

\begin{definition}[An exact admitted-test translation]
A family $g_a:D_{Fa}\to E_a$ is exact when
\begin{equation}\label{eq:chu}
 f_a(x)\,S_{Fa}\,d\quad\Longleftrightarrow\quad x\,R_a\,g_a(d).
\end{equation}
For each operation, supply continuation maps $v_u:E_b\to E_a$ and
$\widetilde v_u:D_{Fb}\to D_{Fa}$ satisfying the acted-incidence identities
\[
 u(x)R_b e\Longleftrightarrow xR_av_u(e),\qquad
 \widetilde u(y)S_{Fb}d\Longleftrightarrow yS_{Fa}\widetilde v_u(d).
\]
Also require
\begin{equation}\label{eq:test-nat}
 g_a\widetilde v_u=v_ug_b.
\end{equation}
These are equations between labelled tests. An extensional version may instead
identify tests with their truth predicates, but that is a further, explicit
forgetting step.
\end{definition}

The direction of $g_a$ is forced: a target test is evaluated on a translated
source state by pulling the test back. A compositional operational translation
can prove this identity for tests represented by translated native probes.
It does not show that every target probe has such a source representation.
Even when the Boolean predicate $d\circ f_a$ is mathematically definable,
it need not be an \emph{admitted} source observation.

\begin{lemma}[Exact representation test]\label{lem:test-representation}
For fixed $a$ and given $f_a$, an exact $g_a$ exists precisely when every
pulled-back target column $f_a^{-1}[C_d]$ is an available source column $C_e$.
For finite test sets this is a directly decidable condition. This columnwise
criterion does not by itself guarantee the coherent labelled choices in
\eqref{eq:test-nat}.
\end{lemma}
\begin{proof}
Equation \eqref{eq:chu} is precisely $f_a^{-1}[C_d]=C_{g_a(d)}$ for each $d$.
Choosing one representing label for each column proves the converse (with the
usual choice qualification for an arbitrary infinite family). Different labels
can represent the same column without satisfying a specified operation on
labels, which explains the final qualification.
\end{proof}

Replacing ``available column'' by ``closed extent'' is enough for some concept
lattice transport, but is weaker: an intersection of native tests need not be
one admitted native test, and can change the common-witness topology. There is
no reason to strengthen the observation language unnoticed.

\section{The resulting characterization comparison}
For a Boolean context $R\subseteq X\times E$, write
$U'=\{e:(\forall x\in U)\ xRe\}$, $B'=\{x:(\forall e\in B)\ xRe\}$,
and $c_R(U)=U''$. Let $Q_R$ be the quotient by equality of rows,
$\mathcal L_R$ the lattice of closed extents, and $K_R$ the object Dowker
complex. Its nonempty faces are finite sets $U\subseteq X$ with $U'\ne\varnothing$.
For the supported concept comparison below assume $X,E$ finite, as in the
finite part of \cite{Common}. The row and lattice statements are not finite
in essence.

\begin{theorem}[Two separate gates give the coherent bridge]\label{thm:bridge}
Suppose an operational comparison has an interpreted state map satisfying
\eqref{eq:state-nat}, and an exact admitted-test translation satisfying
\eqref{eq:chu}--\eqref{eq:test-nat}. Then for every type it induces:
\begin{enumerate}
\item A row map $\bar f:Q_R\to Q_S$, $[x]\mapsto[f(x)]$, intertwining
all translated admitted operations.
\item The canonical complete-join map and its right adjoint
\[
 f_!:\mathcal L_R\rightleftarrows\mathcal L_S:f^*,\qquad
 f_!(A)=c_S(f[A]),\quad f^*(B)=f^{-1}[B].
\]
It carries marked object concepts to their translated concepts and
intertwines the lifted operations.
\item The maximal partial simplicial comparison
\[
 K_R\longleftarrow K_f\xrightarrow{f}K_S,
 \quad K_f=\{\varnothing\}\cup
 \{\varnothing\ne U\in K_R:U'\cap\im(g)\ne\varnothing\}.
\]
It agrees under closure with the corresponding partial map on supported
concepts, retaining exactly this domain.
\end{enumerate}
For a second such comparison $(p,h)$, the composite is $(pf,gh)$.
The row maps and lattice maps compose, and the partial domain satisfies
\[
 K_{pf}=\{U\in K_f:f[U]\in K_p\}.
\]
Operation-translation squares commute in the partial-map category, including
their domains. No row, lattice, homotopy or native-cell equivalence is implied.
\end{theorem}
\begin{proof}
Equation \eqref{eq:chu} says exactly that $(f,g)$ is a Chu morphism. Equal
source rows therefore have equal target rows. Equation \eqref{eq:state-nat}
proves equivariance of the row map.

For a target closed extent $B=D_0'$, one has
$f^{-1}[B]=(g[D_0])'$, which is closed. Hence inverse image restricts to
closed extents and
\[
 c_S(f[A])\subseteq B\ \Longleftrightarrow\ A\subseteq f^{-1}[B]
\]
proves the adjunction. It also gives
$c_S(f[c_R(U)])=c_S(f[U])$ and the marked-generator identity
$f_!c_R(\{x\})=c_S(\{f(x)\})$. These determine the join map uniquely.
Composition follows from the closure identity.

For a nonempty finite $U$, its target image is a face precisely when
\[
 (f[U])'=g^{-1}[U']\ne\varnothing,
\]
which is the displayed domain criterion. Closing $U$ does not change $U'$;
therefore the same condition governs supported concepts and makes the closure
comparison commute with the span. For $(p,h)$ the criterion becomes
$U'\cap\im(gh)\ne\varnothing$, equivalently
$g^{-1}[U']\cap\im(h)\ne\varnothing$, proving exact partial composition.
The two Chu composites in an operation-translation square are equal by
\eqref{eq:state-nat} and \eqref{eq:test-nat}; applying these formulas proves
all remaining assertions. These are the existing algebraic and partial
geometric calculations of \cite{Common}, with the additional source-level
hypotheses explicitly exhibited.
\end{proof}

\begin{proposition}[Exact row adequacy]
The map $\bar f$ in Theorem~\ref{thm:bridge} is injective exactly when the
pulled-back target tests distinguish every pair of distinct source rows.
It is surjective exactly when every target row is represented by some
translated source state. Thus a strong contextad morphism can induce a proper
quotient, a proper inclusion, or neither on observable states.
\end{proposition}
\begin{proof}
Two source row classes have equal images precisely when all target tests have
equal responses after translation, equivalently all tests in $\im(g)$ agree
on them. Surjectivity is exactly the definition of target row representation.
\end{proof}

\subsection{Enrichment is oriented differently in the three views}
If the underlying state set is fixed and $E_0\subseteq E_1$ adds observations,
the identity state map with test inclusion is a Chu morphism
$R_1\to R_0$, not $R_0\to R_1$. It gives a surjection
$Q_1\twoheadrightarrow Q_0$ and a join-preserving closure map
$\mathcal L_1\to\mathcal L_0$. By contrast, witnesses are added to the
Dowker construction, so $K_0\hookrightarrow K_1$. The Chu partial span is
\[
 K_1\longleftarrow K_0\xrightarrow{\id}K_0.
\]
This is not a paradox, nor a covariance convention to be repaired. The
characterizations use different native morphism classes. The universal-test
circle-to-disk example in \cite{Common} already shows that unchanged row
states and unchanged full concept lattice do not force unchanged topology.

\begin{example}[An added operation forces a new state characterization]
Let $X=\{0,1,2\}$, let the only old operation be the identity, and read
$h(0)=h(1)=0$, $h(2)=1$. The old observational quotient has two states.
Add the idempotent operation $u(0)=0$, $u(1)=u(2)=2$. Its continuation
observation $h\circ u$ distinguishes $0$ and $1$, so the closed test family
$\{h,h\circ u\}$ has three distinct rows. The operation $u$ cannot descend
to the old quotient, because equivalent old states $0,1$ have inequivalent
images $0,2$. The canonical comparison is from the new three-state quotient
to the old two-state quotient. There is no reverse comparison commuting with
the two quotient maps from $X$. This is a direct deterministic obstruction,
not an assertion that every addition of an operation is a contextad morphism.
\end{example}

\section{A concrete positive contextad change}
\begin{example}[Refining an available reader environment]\label{ex:reader}
Let $r:E'\twoheadrightarrow E$ be a surjection of nonempty sets. The reader
comonads have $G_E(A)=E\times A$ and $G_{E'}(A)=E'\times A$. Their coKleisli
composition is
\[
 (g\star f)(e,a)=g(e,f(e,a)).
\]
The lineator $r\times\id_A:G_{E'}A\to G_EA$ defines a contextad morphism
in the direction from the $E$-reader doctrine to the $E'$-reader doctrine,
acting on bodies by
\[
 T_r(f)(e',a)=f(r(e'),a).
\]
Indeed,
\[
 T_r(g\star f)(e',a)
 =g(r(e'),f(r(e'),a))
 =(T_r(g)\star T_r(f))(e',a),
\]
and identities are preserved. This is the comonad instance of
\cite[Example 5.33]{CM2}. Notice the reversal: $r$ runs from the richer
environment to the coarse one, while the operational embedding runs from
coarse reader programs to rich reader programs.

The functor is faithful on each hom set: surjectivity of $r$ means unequal
source functions remain unequal after pullback. Its image consists exactly
of bodies constant on every $r$-fibre. It is not full when $r$ has a nontrivial fibre: choose input set $A=1$
and a two-element output set. A richer body can distinguish points in that
fibre, while every translated source body is constant on it.

For $E=\{*\}$ and $E'=\{0,1\}$, old bodies ignore the environment bit.
The rich probe $d(e')=e'$ does not descend through $r$: there is no
$h:E\to\{0,1\}$ with $d=h\circ r$. Thus the operation embedding is
completely compositional while new environment observations need not be
expressible using the old environment. Restricting to fibre-constant probes
makes downward observation descent exact; adding $d$ genuinely refines the
observational quotient. Here $r:E'\to E$ compares environment values; it
is not a canonical state map $E\to E'$ in the direction of $T_r$. To apply
Theorem~\ref{thm:bridge}, an actual state interpretation still has to be
specified. This example proves the separate obstruction to downward test
descent, not a reversal of the Chu test variance.
\end{example}

\section{Display spans and the exact composition obstruction}
Let $C$ have a pullback-stable, composition-stable class of display maps.
The display contextad has grades $p:X\to A$ and extension $e_Ap=X$.
A contextual arrow is the span $A\xleftarrow{p}X\xrightarrow{f}B$.
Composition uses the actual pullback $X\times_BY$.
If $F:C\to C'$ sends display maps to display maps and preserves the required
pullbacks, it induces a contextad morphism and therefore a strong span
comparison. This is \cite[Example 5.35]{CM2}; the strong compositor is the
inverse of the canonical isomorphism
\begin{equation}\label{eq:pullback-map}
 \alpha:F(X\times_BY)\longrightarrow FX\times_{FB}FY.
\end{equation}

Without pullback preservation the displayed map still exists. Its direction
is from the image of the composite to the composite of the images. Under the
usual forward-apex-map convention for span cells, it is an \emph{oplax}
compositor. It is not the lax comparison
$Tg\circ Tf\to T(g\circ f)$ used in the relative manuscript. Choosing an
inverse or reversing this arrow without justification changes the claim.
The canonical maps satisfy oplax unit and associativity coherences: both
routes from an image of an iterated pullback to the target iterated pullback
have the same maps to each factor, and uniqueness in the target pullback
forces equality, with the usual associators understood.

\begin{example}[Forgetting a symmetry destroys relative correlation]
Let $G=\mathbb Z/2$ and let $F:G\text{-}\Set\to\Set$ take orbits.
Use all maps as displays. Let $B=1$, and let $X=Y=G$ with their regular
left action. Then $FX=FY=1$, while $F(X\times Y)$ has two elements:
\[
 \{(0,0),(1,1)\},\qquad \{(0,1),(1,0)\}.
\]
Thus \eqref{eq:pullback-map} is the map $\{\text{equal},\text{different}\}\to1$.
It is not invertible. Each factor has a single orbit, but the composite
retains a relative bit. The invariant joint test $x=y$ distinguishes its
classes. No test on the singleton composite of the separate orbit quotients
can represent this distinction.

For finite groups, more generally, the fibre of the map over an orbit pair
$([x],[y])$ in $FX\times_{FB}FY$ is the set of diagonal $G$-orbits of
compatible representative pairs. It records relative alignment, not merely
two independent marginal orbit labels. The two-element calculation above
is enough to disprove strong compositionality. It does not disprove an
oplax semantics retaining the displayed comparison, and it does not claim
to be a new obstruction to orbit functors preserving limits.
\end{example}

\subsection{How this fits the previous wiring theorem}
The wiring manuscript \cite{Wiring} uses a concrete boundary-reduction square
and the canonical map
\[
 X\times_BY\longrightarrow Q_X\times_{\bar B}Q_Y.
\]
Its image measures jointly realizable reduced pairs. When this map is not
surjective, the full reduced pullback has phantom compatibility. When it is
surjective but not injective, distinctions of joint realization can still be
lost. The orbit example exhibits the second kind and shows that even an
honest functor on bases need not induce a \emph{strong} display-span map.
An arbitrary family of state quotients need not assemble into a base functor
at all, so it is weaker input still.

The orbit example also explains why admissible exterior probes matter. If
the response language includes relative equality, the singleton cannot be
an exact behavioral replacement. If that test is deliberately excluded,
the singleton may be adequate for the selected responses while remaining
inadequate for native witness recovery. The choice must be explicit.

\section{What still has to be proved in an actual application}
The contextad framework contributes typed operational syntax, dependent
reindexing, coherent composition, endpoint squares, and structured morphisms
between such packages. It prevents treating a comparison on objects as a
comparison of whole operational languages. The Chu bridge then gives the
characterization comparison only after its separate observational premises
are met.

For a proposed application the following obligations remain substantive.
\begin{enumerate}
\item Specify the native contextad or other operational doctrine, not only its
objects or a finite response table; give its grades, extension, units,
accumulation and needed coherence.
\item Construct the language-change morphism, including its lineator,
cartesian transport, unit and multiplication comparisons. If only an oplax
or lax map exists, retain that exact direction and its cells.
\item Give the state interpretation and the admitted readouts. Prove which
continuations are permitted, including access, nonrepeatability, resources
and shared state.
\item Represent every translated test by an admitted source test, prove exact
response compatibility and continuation naturality, or explicitly restrict
the comparison to the tests for which this holds.
\item Separately establish joint lifting and external separation for wiring.
Contextad coherence does not make a quotient preserve pullbacks or make all
external observations separating.
\item If native realization or witness preservation is intended, supply
native realizations and cartesian witness comparisons, or an appropriate
fully faithful/essentially surjective categorical result. Boolean row
adequacy does not recover discarded certificate fibres or endpoint squares.
\end{enumerate}

This is a genuine, conditional bridge to the user's earlier contextad work.
It does not assign a contextad structure to every observation change, assert
that all characterizations have one variance, or claim that a behavioral
quotient reconstructs the native contextual system.

\begin{thebibliography}{9}
\bibitem{CM1} M. Capucci and D. J. Myers.
\emph{Contextads as Wreaths; Kleisli, Para, and Span Constructions as Wreath
Products}. arXiv:2410.21889v1 (2024), 84 pages.
\url{https://arxiv.org/pdf/2410.21889v1}.

\bibitem{CM2} M. Capucci and D. J. Myers.
\emph{Contextful Arrows via Wreath Products}.
arXiv:2410.21889v2, 3 November 2026, 120 pages.
Especially Definition 5.25, Examples 5.33--5.35, Definition 5.36,
equations (5.6.20)--(5.6.21), and the qualifications in Section 5.4 and
Remark 5.20. \url{https://arxiv.org/pdf/2410.21889v2}.

\bibitem{Relative} Prior project manuscript.
\emph{One relative construction for context content and effects},
1 November 2026. Sections 2, 4 and 9: base-changing operational comparisons,
marked presentation recovery and controlled forgetting.

\bibitem{Common} Prior project manuscript.
\emph{One operational construction for residual automata, formal concepts,
and topology}, 4 November 2026. Sections 3 and 5--8: exact Chu lifting,
maximal partial domains and variance distinctions.

\bibitem{Wiring} Prior project manuscript.
\emph{Contextual Reduction and Wiring for Open Learners: Behavioral images,
native optimization, and structural characterizations}, 5 November 2026.
Source \nolinkurl{contextual_reduction_wiring.tex}.
Typed contextual reduction, canonical comparison, joint lifting, and
native witness obligations.
\end{thebibliography}
\end{document}
